\section{Schmidt-number geometry and entanglement witnesses}
\label{sec:app_positive_not_cp_schmidt_number_geometry}

This section supplies the convexity, compactness, and separation results used
in the witness criterion of Theorem~\ref{thm:not_has_schmidt_number_le_iff_exists_witness}.

\begin{theorem}[Basic properties of bounded Schmidt number]
    \label{thm:has_schmidt_number_le_basic}
    \lean{Matrix.HasSchmidtNumberLE.posSemidef, Matrix.HasSchmidtNumberLE.add,
        Matrix.HasSchmidtNumberLE.mono, Matrix.HasSchmidtNumberLE.smul}
    \leanok
    \uses{def:has_schmidt_number_le}
    A matrix of Schmidt number at most $n$ is positive semidefinite. States of
    Schmidt number at most $n$ are closed under addition and under multiplication by a
    non-negative real scalar, and a bound on the Schmidt number relaxes to any larger
    bound.
\end{theorem}

\begin{proof}
    \leanok
    Each summand $\ket{\psi_i}\bra{\psi_i}$ is a rank-one positive semidefinite
    matrix, and a finite sum of positive semidefinite matrices is positive
    semidefinite. Concatenating two pure-state families realizes a sum as a single
    finite sum of pure-state projectors of Schmidt rank at most $n$, and a pure
    summand of Schmidt rank at most $n$ also has Schmidt rank at most any larger
    bound. Multiplication by a non-negative real $a$ rescales each pure summand
    $\psi_i$ to $\sqrt a\,\psi_i$, which leaves its Schmidt rank unchanged and
    reproduces $a$ times the original state.
\end{proof}

\begin{theorem}[Convexity of the Schmidt-number set]
    \label{thm:convex_setOf_has_schmidt_number_le}
    \lean{Matrix.convex_setOf_hasSchmidtNumberLE}
    \leanok
    \uses{def:has_schmidt_number_le}
    For each $n$ the set $S_n$ of bipartite states of Schmidt number at most $n$ is
    convex.
\end{theorem}

\begin{proof}
    \leanok
    A convex combination $a x + b y$ with $a,b\ge 0$ is a sum of two non-negatively
    scaled states of Schmidt number at most $n$, and that bound is preserved under
    non-negative scaling and under addition; rescaling a vector by a non-negative
    scalar scales its coefficient matrix and so does not increase the Schmidt rank.
\end{proof}

\begin{theorem}[Compactness of the Schmidt-number state set]
    \label{thm:compact_setOf_has_schmidt_number_le_trace_one}
    \lean{Matrix.isCompact_setOf_hasSchmidtNumberLE_trace_one}
    \leanok
    \uses{def:has_schmidt_number_le, def:schmidt_rank,
        thm:convex_setOf_has_schmidt_number_le}
    For each $n$ the set $S_n$ of trace-one bipartite states of Schmidt number at most
    $n$ is compact.
\end{theorem}

\begin{proof}
    \leanok
    A trace-one state of Schmidt number at most $n$ is a convex combination of pure
    states $\ket{\psi}\bra{\psi}$ whose vector $\psi$ has unit norm and Schmidt rank
    at most $n$, so $S_n$ is the convex hull of the set $P_n$ of those pure states.
    The set $P_n$ is the image, under the continuous projector map
    $\psi\mapsto\ket{\psi}\bra{\psi}$, of the unit vectors of Schmidt rank at most
    $n$. In the Euclidean norm those unit vectors form the unit sphere, which is
    compact in finite dimension, intersected with the closed set on which the Schmidt
    rank is at most $n$ (the coefficient matrix depends continuously on the vector and
    the matrices of rank at most $n$ form a closed set); hence $P_n$ is compact. The
    trace of $\ket{\psi}\bra{\psi}$ is the squared Euclidean norm of $\psi$, so the
    trace-one constraint is exactly the unit-norm constraint and the weights of the
    convex decomposition are the squared norms, summing to the trace. The convex hull
    of a compact set is compact in a finite-dimensional real normed space.
\end{proof}

\begin{theorem}[Trace-form representation of a real-linear functional]
    \label{thm:exists_isHermitian_trace_form_re}
    \lean{Matrix.exists_isHermitian_trace_form_re}
    \leanok
    Every real-linear functional $g$ on the space of square complex matrices is the
    trace form of a Hermitian matrix: there is a Hermitian $H_0$ with
    $g(X)=\re\tr(X H_0)$ for every Hermitian $X$.
\end{theorem}

\begin{proof}
    \leanok
    The real bilinear pairing $(X,H)\mapsto\re\tr(X H)$ is nondegenerate,
    because on the diagonal pair $(H,H^{\dagger})$ it equals the squared Frobenius norm
    $\re\tr(H H^{\dagger})=\sum_{i,j}\lvert H_{ij}\rvert^2$, which is positive
    unless $H=0$. Hence the map sending $H$ to the functional
    $X\mapsto\re\tr(X H)$ is an injective real-linear map from the matrix
    space to its real dual. These two spaces have the same finite real dimension, so the
    map is also surjective, and every $g$ is represented by some $H$. Replacing $H$ by
    its Hermitian part $\tfrac12(H+H^{\dagger})$ keeps the value on Hermitian inputs,
    since $\tr(X H^{\dagger})=\overline{\tr(X H)}$ when $X$ is Hermitian, and produces a
    Hermitian representing matrix.
\end{proof}

\begin{theorem}[Entanglement witness from a separating hyperplane]
    \label{thm:exists_entanglement_witness}
    \lean{Matrix.exists_isHermitian_witness}
    \leanok
    \uses{def:has_schmidt_number_le, def:schmidt_rank,
        thm:convex_setOf_has_schmidt_number_le,
        thm:compact_setOf_has_schmidt_number_le_trace_one,
        thm:exists_isHermitian_trace_form_re}
    Let $\rho$ be a trace-one Hermitian bipartite state whose Schmidt number exceeds $n$.
    Then there is a Hermitian operator $W$ such that, for every $\psi$ of Schmidt rank
    at most $n$,
    \begin{align}
        \re\tr(W\rho)            & <0,
        \notag                            \\
        \re\bra{\psi}W\ket{\psi} & \ge 0.
        \notag
    \end{align}
    This is the only-if direction of Wolf's Proposition~3.3
    \cite[Chapter~3, Proposition~3.3]{Wolf2012Quantum}: a state of Schmidt number larger
    than $n$ is detected by an entanglement witness for $S_n$.
\end{theorem}

\begin{proof}
    \leanok
    The set $S_n$ of trace-one states of Schmidt number at most $n$ is convex and
    compact, and $\rho\notin S_n$. The geometric Hahn--Banach theorem separates the
    closed convex set $\{\rho\}$ from the compact convex set $S_n$ by a continuous
    real-linear functional $f$ and a constant $c$ with $f(\rho)<c<f(\sigma)$ for every
    $\sigma\in S_n$. Representing $f$ as the trace form of a Hermitian $H_0$ and setting
    $W=H_0-c\,\Id$ gives $\re\tr(W\rho)=f(\rho)-c<0$. For a vector $\psi$
    of Schmidt rank at most $n$ the projector $\ket{\psi}\bra{\psi}$, after
    normalization by $\lVert\psi\rVert^2$, is a trace-one state of $S_n$, so its value
    under $f$ is at least $c$; multiplying back by the squared norm gives
    $\re\bra{\psi}W\ket{\psi}\ge 0$, with the zero vector giving the value
    $0$ directly.
\end{proof}

\section{Entanglement witnesses and Choi trace pairing}
\label{sec:app_positive_not_cp_witness_choi}

The following results identify witnesses with Choi matrices of $n$-positive
maps and supply the trace identity used in
Theorem~\ref{thm:exists_isNPositiveMap_detects}.

\begin{theorem}[Surjectivity of the Choi correspondence]
    \label{thm:exists_choiMatrix_eq}
    \lean{ChoiJamiolkowski.exists_choiMatrix_eq}
    \leanok
    \uses{thm:cp_iff_choi_posSemidef}
    Assume $D>0$. Every bipartite matrix $W$ on $\C^D\otimes\C^D$ is the Choi matrix
    $\tau_T=(T\otimes\Id)(\ket{\Omega}\bra{\Omega})$ of some linear map
    $T\colon\MN{D}\to\MN{D}$.
\end{theorem}

\begin{proof}
    \leanok
    The Choi correspondence $T\mapsto\tau_T$ is complex-linear and injective. Its domain,
    the space of linear maps on $\MN{D}$, and its codomain, the bipartite matrix space on
    $\C^D\otimes\C^D$, both have complex dimension $D^4$, so the injective map is also
    surjective. Hence every $W$ is a Choi matrix.
\end{proof}

\begin{theorem}[A witness is the Choi matrix of an $n$-positive map]
    \label{thm:exists_isNPositiveMap_choiMatrix_eq}
    \lean{ChoiJamiolkowski.exists_isNPositiveMap_choiMatrix_eq_of_witness}
    \leanok
    \uses{def:k_positive_map, def:schmidt_rank, thm:exists_choiMatrix_eq,
        thm:schmidt_rank_choi_test_iff, thm:exists_entanglement_witness}
    Assume $D>0$. Let $W$ be a Hermitian operator on $\C^D\otimes\C^D$ with
    $\re\bra{\psi}W\ket{\psi}\ge 0$ for every $\psi$ of Schmidt rank at most $n$.
    Then $W$ is the Choi matrix of an $n$-positive map $T$. This is the
    Choi--Jamio\l kowski translation of the entanglement witness of Wolf's Proposition~3.3
    into the $n$-positive map of Wolf's Proposition~3.4
    \cite[Chapter~3, Proposition~3.4]{Wolf2012Quantum}.
\end{theorem}

\begin{proof}
    \leanok
    By the surjectivity of the Choi correspondence (\ref{thm:exists_choiMatrix_eq}) there
    is a linear map $T$ with $\tau_T=W$. Because $W$ is Hermitian, the Choi quadratic form
    $\bra{\psi}W\ket{\psi}$ is real, and it equals the witness expectation
    $\tr(W\ket{\psi}\bra{\psi})$; the witness condition makes it non-negative on every
    vector of Schmidt rank at most $n$. This is exactly the Schmidt-rank Choi criterion
    (\ref{thm:schmidt_rank_choi_test_iff}) for $n$-positivity of $T$.
\end{proof}

\begin{theorem}[Choi trace pairing through the trace-pairing adjoint]
    \label{thm:trace_choiMatrix_omega_quadratic_adjoint}
    \lean{ChoiJamiolkowski.trace_choiMatrix_mul_eq_omegaVec_quadraticForm_traceAdjointMap}
    \leanok
    \uses{def:trace_pairing_adjoint, thm:trace_pairing_adjoint_ampliation}
    For a linear map $T\colon\MN{D}\to\MN{D}$ with Choi matrix $\tau_T$ and any bipartite
    matrix $\rho$ on $\C^D\otimes\C^D$,
    \begin{align}
        \tr(\tau_T\,\rho)
         & =\bra{\Omega}(T^{*}\otimes\Id)(\rho)\ket{\Omega},
        \notag
    \end{align}
    where $T^{*}$ is the trace-pairing adjoint of $T$.
\end{theorem}

\begin{proof}
    \leanok
    Since $\tau_T=(T\otimes\Id)(\ket{\Omega}\bra{\Omega})$, cyclicity of the trace gives
    $\tr(\tau_T\rho)=\tr(\rho\,(T\otimes\Id)(\ket{\Omega}\bra{\Omega}))$. The
    $\Id$-ampliation of the trace-pairing adjoint is the trace-pairing adjoint of the
    ampliation, so moving $T$ across the trace pairing replaces
    $(T\otimes\Id)$ by $(T^{*}\otimes\Id)$ acting on $\rho$ and turns the pairing against
    $\ket{\Omega}\bra{\Omega}$ into the quadratic form $\bra{\Omega} \cdot\,\ket{\Omega}$.
\end{proof}
